\chapter*{Abstract}
\section*{}
Cross-sectional stock-return prediction with neural networks commonly optimises Mean Squared Error (MSE) or a close relative, yet the downstream decision --- a long-short portfolio --- depends on the cross-sectional rank of the predictions, not on their calibrated magnitudes. Under this mismatch, a quadratic loss gives disproportionate training weight to heavy-tailed outliers, and the resulting model trades ranking quality for calibration quality that the portfolio does not use. This report asks whether a hybrid loss that combines a directional-accuracy term with a robust magnitude term can outperform both traditional regression losses and pure directional losses under a controlled evaluation protocol.

The study uses a CRSP-style US equity monthly panel with a static protocol: training window \texttt{1990-01} to \texttt{1994-12} (60 monthly cross-sections) and main out-of-sample test window \texttt{1995-01} to \texttt{1996-12} (24 monthly cross-sections). The model is a fixed multi-layer perceptron (MLP) with hidden widths $[64, 32, 16]$, ReLU activations, and dropout $0.2$, consuming 15 model-input columns (see Chapter 4 \S4.4.1 and Appendix B Table B.3). Training uses batch size $1024$ and 20 epochs of Adam; the trained model is then evaluated month-by-month on a top/bottom-10\% signal-tilted long-short portfolio with a 5\% per-name cap. Performance is reported as annualised Sharpe ($\sqrt{12}$-scaled monthly mean/std), cumulative 24-month return, and, for multi-seed phases, the coefficient of variation $\mathrm{CV} = \sigma_S / |\mu_S|$ across three random seeds.

Four comparison phases plus one normalisation probe are run. Within each comparison table, the loss function is the only varying factor; cross-phase comparisons differ in additional dimensions (\S3.8). A Phase 1 single-seed baseline compares seven loss families (MSE, MedSE, MADL, GMADL, IMADL, and two hybrid multiplicative variants). A Phase 2 sweep scans nine parameterised hybrid variants (five additive A-series and four multiplicative M-series). Phase 3 refines the M2-robust loss with a $\gamma$ parameter and evaluates five $\gamma$ settings across three seeds, then extends the comparison to IMADL-m2 $\alpha$, IMADL-GMADL $\beta$, and adaptive-$\lambda$ sweeps under the same multi-seed protocol. Finally, a Phase 4 normalisation probe tests loss-component normalisation on the top three candidates as a diagnostic follow-up.

The multi-seed evidence identifies \texttt{m2\_robust\_gamma07} (multiplicative hybrid with robustness parameter $\gamma = 0.07$) as the best-supported choice. It achieves mean Sharpe $0.9156$, CV $0.1808$, and mean cumulative return $+27.99\%$ across three seeds; no seed produces a negative Sharpe in this configuration, and the result is approximately flat under loss-component normalisation (normalised mean Sharpe $0.9112$). The high-return alternative \texttt{m2\_robust\_gamma10} reaches a mean Sharpe of $1.0043$ but has CV $0.5613$ --- three times that of $\gamma$07 --- and degrades to mean Sharpe $0.4072$ under normalisation; it is only appropriate when higher seed sensitivity is acceptable. The stable fallback \texttt{imadl\_m2\_alpha06}, drawn from an independent parameterisation, achieves mean Sharpe $0.6895$ with CV $0.2443$ and the largest mean cumulative return of the three candidates at $+30.42\%$; it is selected by the stability/return trade-off rather than by Sharpe dominance. Traditional regression losses (MSE, MedSE) and pure directional losses (MADL, GMADL, IMADL) do not produce economically meaningful positive Sharpes in the seed-42 same-window baseline, and the absolute-loss family additionally exhibits very negative point-prediction R\textsuperscript{2} values while retaining competitive ranking performance --- a dissociation between calibration and portfolio quality that the report makes explicit.

Three boundaries constrain the conclusions. First, the evaluation uses a single static 24-month window; the rolling-window extension is listed as future work. Second, multi-seed evaluations use three seeds, which is sufficient to differentiate CV at the order-of-magnitude level but not to bound CV to the second decimal. Third, the loss-component normalisation probe's scale ratios are diagnostics-estimated rather than measured by a per-component logger, which is flagged as a priority instrumentation task. Within these boundaries the report provides a verifiable recommendation --- \texttt{m2\_robust\_gamma07} as the primary loss, with explicit high-return and stable-fallback alternatives --- together with a source-of-truth evidence map that ties every reported number to a named artifact file.

\vspace{5cm}

\noindent \textbf{KEY WORDS}: loss function design; robust regression; MADL; GMADL; cross-sectional stock-return prediction; long-short portfolio; multi-layer perceptron; multi-seed robustness; coefficient of variation; CRSP
